\section{From Threats to Controls: Practical Security Implications}
\label{sec:discussion}

\subsection{Control Placement Should Follow Failure Origin}
\label{subsec:implications}

The threat analysis in Section~\ref{sec:security} suggests a simple but important lesson: defenses should be attached where failures are \emph{introduced}, not only where they are eventually observed. A CUA that fails at runtime may require intervention during creation, deployment, operation, maintenance, or several of these stages at once.

\textbf{Creation-stage controls} should shape what the model is allowed to learn and emit. Safe action abstractions, constrained high-impact operations, and benchmark-based release gating reduce the chance that unsafe behavior is built into the agent from the outset~\cite{kuntz2025harm,tian2025measuring,miculicich2025veriguard}. This is also the stage at which threat models should be made explicit rather than delegated to downstream operators.

\textbf{Deployment-stage controls} should govern authority and mediation. Least privilege, provenance-aware interfaces, sandboxing, and intent-aware access control are essential because deployment determines which tools, accounts, and execution channels are reachable in practice~\cite{gong2025secure,meng2025cellmate,foerster2026camels}. In open systems, deployment is not a thin wrapper around a model; it is part of the security perimeter.

\textbf{Operation-stage controls} should contain unsafe trajectories before they become high-impact outcomes. Runtime verification, human confirmation for consequential actions, misalignment detection, and defensive treatment of tool outputs all belong here~\cite{kang2025mitigating,ning2026actions,zhang2025browsesafe}. The main question is not only whether an agent can complete a task, but whether it can remain on-task under ambiguity, adversarial inputs, and long-horizon pressure.

\textbf{Maintenance-stage controls} should preserve validity over time. Continuous regression testing, vulnerability management, audit trails, and policy updates are required because interface drift, tool churn, and ecosystem change can reintroduce previously solved problems~\cite{zheng2026riskybench,wang2026asset,liu2026clawkeeper}. For CUAs, secure deployment is therefore not a one-time event.

Under the attribution framework, the control logic becomes more specific. \textbf{Scope overreach} is best contained by clearer task boundaries, data-minimization policy, confirmation thresholds, and memory scoping. \textbf{Objective corruption} is best contained by provenance separation, tool and extension vetting, memory hygiene, retrieval controls, and defenses that prevent external instructions from becoming operative goals. \textbf{Environmental misbinding} is best contained by stronger mediation between perception and execution, least privilege, sandboxing, TOCTOU-aware verification, and auditable action binding.

\subsection{Stakeholder Responsibilities Are Coupled}

\begin{table*}[!t]
\centering
\scriptsize
\caption{Stakeholder responsibilities derived from the joint framework.}
\label{tab:actor_control_matrix}
\renewcommand{\arraystretch}{1.08}
\setlength{\tabcolsep}{3.5pt}
\begin{tabularx}{\textwidth}{L{1.8cm} L{2.8cm} L{2.8cm} L{2.8cm} Y}
\toprule
\textbf{Actor} & \textbf{Creation} & \textbf{Deployment} & \textbf{Operation} & \textbf{Maintenance} \\
\midrule
Developers & Define action abstractions, training objectives, benchmark targets & Expose policy hooks, provenance fields, safe defaults & Support verifier hooks and interpretable traces & Patch regressions and update evaluation criteria \\
Deployment teams & Validate model-package integrity and release gates & Scope privileges, sandbox tools, configure trust boundaries & Enforce confirmation thresholds and environment isolation & Reconfigure permissions and retire unsafe integrations \\
Operators & Define task policy and escalation rules & Approve high-impact workflows and account bindings & Monitor incidents, review risky actions, handle exceptions & Feed failures back into policy and regression suites \\
Ecosystem stewards & Maintain registries, signing, disclosure process & Govern plugin, skill, and MCP distribution & Detect abuse propagation and preserve identity/provenance & Manage coordinated patching, version pinning, ecosystem response \\
\bottomrule
\end{tabularx}
\end{table*}

The same reasoning also changes how responsibilities should be assigned. Developers mainly own creation-stage risks such as unsafe abstractions, training objectives, and latent objective corruption. Deployment teams mainly own environmental-misbinding risks introduced by privilege, tooling, and mediation choices. Operators mainly own operation-stage scope control, escalation, and runtime detection. Ecosystem stewards mainly own maintenance-stage propagation, registry integrity, and coordinated response. These roles are analytically separable but operationally coupled: weak deployment can negate careful model training, and weak maintenance can erase earlier safety gains.

Viewed through the publicly described OpenClaw example, this division of labor becomes concrete. Developers define what actions can be safely exposed; deployment teams decide how messaging channels, tools, and sessions are bound; operators govern approval and escalation for live tasks; and ecosystem stewards manage the provenance and policy discipline of the broader capability-sharing surface. This coupling is especially visible in open ecosystems, where provenance, identity, and version discipline become preconditions for safe operation~\cite{chen2026openclaw,wang2026asset,liu2026clawkeeper}.

\subsection{A Minimal Defense Stack for Open CUA Systems}
\label{subsec:recommendations}

No single mechanism is sufficient. The current literature instead points toward a minimal layered defense stack for realistic CUA deployment:
\begin{itemize}
  \item \textbf{Constrained action interfaces.} High-impact operations should be represented through auditable, semantically narrow interfaces rather than unconstrained low-level action spaces wherever possible~\cite{jones2025systematization,miculicich2025veriguard}.
  \item \textbf{Provenance-aware mediation.} The system should preserve the distinction between user-authored intent, retrieved content, tool outputs, and environment text so that instruction source remains inspectable~\cite{kang2025mitigating,zhang2025browsesafe}.
  \item \textbf{Least privilege and sandboxing.} Tool access, filesystem reach, network scope, and account authority should be narrowed before deployment, not after the first incident~\cite{gong2025secure,meng2025cellmate,foerster2026camels}.
  \item \textbf{Runtime verification and escalation control.} High-impact actions should trigger verification, policy checks, or human confirmation, especially under uncertainty or when plan execution deviates from task scope~\cite{ning2026actions,zhang2026mirrorguard,liu2026clawkeeper}.
  \item \textbf{Continuous evaluation and ecosystem governance.} Regression suites, red teaming, signed registries, and versioned identities are necessary to prevent maintenance-stage drift from becoming a supply-chain or multi-agent security problem~\cite{liao2025redteamcua,zheng2026riskybench,zhang2025mcp,wang2026asset}.
\end{itemize}

Safe CUA deployment depends less on any single alignment technique than on whether the full stack connects architectural choices, authority boundaries, and operational controls in a coherent way.
